
\section{Detailed Experimental Setup} \label{appendix:detailed-experimental-setup}
\subsection{Correct and Counterfeit Data Generation for Correctness Checking}
In Table \ref{tab:dataset-sizes-correctness-checking}, we show statistics about the datasets used for correctness checking. Recall that each dataset consists of 5 correct and 5 counterfeit samples per problem. We also show the average pass@1 score of problems in the dataset at $T=0.6$. A few examples of correct and counterfeit samples are shown in Listings \ref{lst:dataset-humaneval-example}, \ref{lst:dataset-leetcode-example}, and \ref{lst:dataset-odex-example}.

\begin{table}[H]
    \centering
    \caption{Correctness Checking Dataset Sizes}
    \begin{tabular}{cccc}
        \toprule
        \textbf{Dataset} & \textbf{Model} & \textbf{Pass@1} & \textbf{Size} \\
	\midrule
	\multirow{5}{*}{HumanEval}
	& CL-34b & 42.0 & 850 \\
	& CL-7b & 36.9 & 870 \\
	& DS-I-33B & 45.1 & 830 \\
	& StarCoder & 32.9 & 660 \\
	& DS-I-6.7B & 56.7 & 810 \\
        \midrule
	\multirow{2}{*}{LeetCode}
	& DS-I-33B & 49.4 & 460 \\
	& DS-I-6.7B & 37.3 & 360 \\
	\midrule
	\multirow{5}{*}{ODEX}
	& CL-34B & 49.2 & 1070 \\
	& CL-7B & 52.2 & 1190 \\
	& DS-I-33B & 44.3 & 520 \\
	& StarCoder & 46.2 & 1060 \\
	& CL-13B & 50.1 & 1090 \\
    \bottomrule
    \end{tabular}
    \label{tab:dataset-sizes-correctness-checking}
\end{table}

\textbf{HumanEval}: HumanEval \citep{chen2021evaluating} is a dataset of 164 relatively simple natural language to programming problems in Python. We sample $200$ generations at $T=0.6$. We use both the original HumanEval tests and EvalPlus tests, which are more comprehensive \citep{liu2023your}. In order to filter out trivial solutions and keep the task interesting, we only consider a counterfeit sample to be a program with an EvalPlus score of over 10\% and manually inspect the resulting dataset. On manual inspection, we found that EvalPlus tests can sometimes be too strong and filter out correct solutions due to very subtle errors like precision and floating point issues, we consider a solution as correct if it passes all the base tests and at least 95\% of EvalPlus tests. Our manual inspection shows that this is a fairer criteria for assessing program correctness.

\textbf{LeetCode}: LeetCode is a dataset of 130 LeetCode problems used for programming interviews. We sample $200$ generations at $T=0.6$. Unlike in HumanEval, run-time and computational complexity is an important consideration for many LeetCode problems. However, determining whether a program can finish within the time limit can be difficult. Therefore, we consider counterfeit programs to be those that resulted in a ``Correct'' or ``Wrong Answer'' verdict, and remove programs that received a ``Runtime Error'' and ``Time Limit Exceeded''. 

\textbf{ODEX}: ODEX \citep{wang2022execution} is an open-domain, multilingual, execution-based natural language to code generation benchmark. We only use the English subset of 479 problems. Unlike HumanEval and LeetCode, ODEX contains problems using a wide variety of Python library functions such as \texttt{numpy}, \texttt{os}, and \texttt{pandas}. We sample $50$ generations at $T=0.6$. As ODEX does not come with cleanly separated test cases, we consider a program as counterfeit if it can be parsed successfully by \texttt{ast.parse} and have a length of under 500 characters. From manual inspection, some of the problems in ODEX can be quite vague, making it difficult to discern if a solution is correct or counterfeit without seeing the input-output format. To mitigate this, for ODEX only, we include \textit{both the generated program and the assertions that is checked}. This leads to a slightly easier setting than the previous two tasks, but we find that it still poses a significant challenge for models.

\begin{lstlisting}[language=python,caption={Example of counterfeit sample for HumanEval, generated by CL-34B},style=boxed,label={lst:dataset-humaneval-example}]
def sort_even(l: list):
    """This function takes a list l and returns a list l' such that
    l' is identical to l in the odd indicies, while its values at the even indicies are equal
    to the values of the even indicies of l, but sorted.
    >>> sort_even([1, 2, 3])
    [1, 2, 3]
    >>> sort_even([5, 6, 3, 4])
    [3, 6, 5, 4]
    """
    even_nums = l[::2]
    even_nums.sort()
    odd_nums = l[1::2]
    ans = []
    for i in range(len(even_nums)):
        ans.append(even_nums[i])
        ans.append(odd_nums[i])
    return ans
\end{lstlisting}

\begin{lstlisting}[language=python,caption={Example of correct sample for LeetCode, generated by DS-I-33B},style=boxed,label={lst:dataset-leetcode-example}]
"""
You are given a 0-indexed permutation of n integers nums.
A permutation is called semi-ordered if the first number equals 1 and the last number equals n. You can perform the below operation as many times as you want until you make nums a semi-ordered permutation:

Pick two adjacent elements in nums, then swap them.

Return the minimum number of operations to make nums a semi-ordered permutation.
A permutation is a sequence of integers from 1 to n of length n containing each number exactly once.
 
Example 1:

Input: nums = [2,1,4,3]
Output: 2
Explanation: We can make the permutation semi-ordered using these sequence of operations: 
1 - swap i = 0 and j = 1. The permutation becomes [1,2,4,3].
2 - swap i = 2 and j = 3. The permutation becomes [1,2,3,4].
It can be proved that there is no sequence of less than two operations that make nums a semi-ordered permutation. 

Example 2:

Input: nums = [2,4,1,3]
Output: 3
Explanation: We can make the permutation semi-ordered using these sequence of operations:
1 - swap i = 1 and j = 2. The permutation becomes [2,1,4,3].
2 - swap i = 0 and j = 1. The permutation becomes [1,2,4,3].
3 - swap i = 2 and j = 3. The permutation becomes [1,2,3,4].
It can be proved that there is no sequence of less than three operations that make nums a semi-ordered permutation.

Example 3:

Input: nums = [1,3,4,2,5]
Output: 0
Explanation: The permutation is already a semi-ordered permutation.

 
Constraints:

2 <= nums.length == n <= 50
1 <= nums[i] <= 50
nums is a permutation.
"""

def semiOrderedPermutation(nums: List[int]) -> int:
    n = len(nums)
    min_idx = nums.index(1)
    max_idx = nums.index(n)
    
    if min_idx < max_idx:
        return min_idx + n - 1 - max_idx
    else:
        return min_idx + n - 2 - max_idx
\end{lstlisting}

\begin{lstlisting}[language=python,caption={Example of counterfeit sample for ODEX, generated by DS-I-33B},style=boxed,label={lst:dataset-odex-example}]
def f_41861705(split_df, csv_df):
    """Create new DataFrame object by merging columns "key" of  dataframes `split_df` and `csv_df` and rename the columns from dataframes `split_df` and `csv_df` with suffix `_left` and `_right` respectively
    """
    return split_df.merge(csv_df, how='left', left_on='key', right_on='key', suffixes=('_left', '_right'))

import pandas as pd

def check(candidate):
    split_df = pd.DataFrame({'key': ['foo', 'bar'], 'value': [1, 2]})
    csv_df = pd.DataFrame({'key': ['foo', 'baz'], 'value': [3, 4]})
    result = pd.DataFrame({'key': ['foo'], 'value_left': [1],'value_right': [3]})
    assert all(candidate(csv_df, split_df) == result)

check(f_41861705)
\end{lstlisting}

\subsection{Data Generation for Execution Prediction}
We perform code execution experiments on HumanEval and LeetCode programs. The inputs and outputs for these datasets are primitive Python objects (mostly \texttt{int}, \texttt{str}, \texttt{bool}, \texttt{list}). While it is possible, we do not evaluate execution for ODEX because many of the programs involve file modifications and cannot easily be represented. For each dataset and data-generating model, we use the same set of programs used in the correctness checking experiment for consistency. As of today, we cannot expect a language model to follow the execution of arbitrary Python programs. Therefore, we ensure that the execution samples in our benchmark are reasonable by applying a filter following the setup in \citep{gu2024cruxeval}. One key difference from their work is that instead of using arbitrary programs, the programs we use here are seeded from a natural language specification and are semantically meaningful. This allows us to analyze how models behave differently when asked to reason about correct and counterfeit programs.

We create our dataset of samples to evaluate code execution as follows: first, we take the programs generated for the correctness checking dataset. The docstring containing the problem statement is stripped away to force the model to use the provided code. Second, we run the program on the tests provided in the original problem statement and examples, which are generally simple and concise to create a large set of model-generated programs, inputs, and outputs. Third, we apply a compile-time and runtime based filter using Python bytecode to remove programs that are too long, require complex arithmetic/floating point operations, and have too many steps in the execution. The final step is a manual inspection of programs, inputs, and outputs passing the filter to ensure that they seem reasonable. The resulting dataset sizes are shown in Table \ref{tab:dataset-sizes-execution}, and examples are shown in Listings \ref{lst:dataset-humaneval-execution-example}, \ref{lst:dataset-leetcode-execution-example}.

\begin{table}[H]
    \centering
    \caption{Execution Dataset Sizes}
    \begin{tabular}{ccc}
        \toprule
        \textbf{Dataset} & \textbf{Model} & \textbf{Dataset Size} \\
	\midrule
	\multirow{5}{*}{HumanEval}
	& CL-34B & 1406 \\
	& CL-7B & 1528 \\
	& DS-I-33B & 1964 \\
	& StarCoder & 1622 \\
	& DS-I-6.7B & 1917 \\
        \midrule
	\multirow{2}{*}{LeetCode}
	& DS-I-33B & 845 \\
	& DS-I-6.7B & 694 \\
    \bottomrule
    \end{tabular}
    \label{tab:dataset-sizes-execution}
\end{table}

\begin{lstlisting}[language=python,caption={Example of HumanEval execution prediction example, generated by StarCoder},style=boxed,label={lst:dataset-humaneval-execution-example}]
from typing import List


def string_xor(a: str, b: str) -> str:
    assert len(a) == len(b)
    res = ""
    for i in range(len(a)):
        if a[i] == b[i]:
            res += "0"
        else:
            res += "1"
    return res
assert string_xor('1', '1') == ??
# Answer: '0'
\end{lstlisting}

\begin{lstlisting}[language=python,caption={Example of LeetCode execution prediction example, generated by DS-I-6.7B},style=boxed,label={lst:dataset-leetcode-execution-example}]
def relocateMarbles(nums: List[int], moveFrom: List[int], moveTo: List[int]) -> List[int]:
    # Create a dictionary to store the number of marbles at each position
    marbles = {}
    for num in nums:
        marbles[num] = marbles.get(num, 0) + 1

    # Apply the moves
    for f, t in zip(moveFrom, moveTo):
        # Remove the marbles at the source position
        count = marbles.pop(f)
        # Add the marbles at the target position
        marbles[t] = marbles.get(t, 0) + count

    # Return the sorted keys of the dictionary
    return sorted(marbles.keys())
assert relocateMarbles(nums = [1, 6, 7, 8], moveFrom = [1, 7, 2], moveTo = [2, 9, 5]) == ??
# Answer: [5, 6, 8, 9]
\end{lstlisting}

